%! Unit 8: Learning from Data — Summative Assessment — Practice Test --- Study Copy (KEY)
\documentclass[10pt]{article}
\usepackage{linalg-article}
\usepackage{linalg-key}   % pulls in -boxes; do NOT also load -boxes
\newcommand{\vv}{\mathbf{v}}
\newcommand{\bb}{\mathbf{b}}
\newcommand{\ww}{\mathbf{w}}
\newcommand{\relu}{\mathrm{ReLU}}

% Part-divider strip
\newcommand{\parthead}[1]{%
  \vspace{6pt}\noindent
  \begin{headlinebox}{burgundy}{\color{white}\bfseries #1}\end{headlinebox}%
  \vspace{4pt}}

\begin{document}

\pageheader{Unit 8: Learning from Data}{Practice Test --- Study Copy --- Answer Key}
\namedateperiod

\vspace{0.10in}
\begin{remindbox}
\small
\textbf{This is a practice test.} It mirrors the real Unit 8 test in length, format,
and the ideas it covers, but the numbers and contexts differ. Work every problem
with no notes, then check your answers against the key.
\end{remindbox}

% ======================================================================
\parthead{Part A --- Vocabulary (8 pts)}
Write the letter of the best description next to each term.

\vspace{4pt}
\noindent
\begin{minipage}[t]{0.40\textwidth}
\begin{enumerate}[leftmargin=2.2em, itemsep=7pt, label=\arabic*.]
  \item ReLU \ans{C}
  \item Neural network layer \ans{A}
  \item Piecewise linear function \ans{F}
  \item Convolution filter \ans{D}
  \item Weight sharing \ans{G}
  \item Gradient descent \ans{E}
  \item Learning rate \ans{B}
  \item Covariance matrix \ans{H}
\end{enumerate}
\end{minipage}%
\hfill
\begin{minipage}[t]{0.57\textwidth}
\small
\begin{description}[leftmargin=1.4em, itemsep=3pt]
  \item[(A)] The building block $\relu(A\vv+\bb)$: multiply by weights, add a bias, then rectify.
  \item[(B)] The step size $s$ in gradient descent --- too small crawls, too big overshoots.
  \item[(C)] The rectifier $\max(0,x)$: keep positive inputs, send negatives to $0$.
  \item[(D)] A small weight pattern slid across the input to detect one feature everywhere.
  \item[(E)] The rule $\ww\leftarrow\ww-s\,\nabla L$: step opposite the slope to shrink the loss.
  \item[(F)] A function built of straight pieces joined at folds (a sum of shifted ReLUs).
  \item[(G)] Reusing one filter at every position, so a giant layer needs only a few weights.
  \item[(H)] The symmetric matrix of variances and covariances; its eigenvectors are the principal components.
\end{description}
\end{minipage}

% ======================================================================
\parthead{Part B --- Multiple Choice (12 pts)}
Circle the best answer.

\begin{enumerate}[leftmargin=1.9em, itemsep=9pt]
  \item $\relu(-3)$ equals
    \hfill (A) $-3$ \quad \textbf{(B) $0$} \ans{$\leftarrow$} \quad (C) $3$ \quad (D) $1$
  \item One neural-network layer computes
    \hfill (A) $A\vv$ \quad \textbf{(B) $\relu(A\vv+\bb)$} \ans{$\leftarrow$} \quad (C) $\vv+\bb$ \quad (D) $A^{-1}\vv$
  \item A piecewise linear function built from $N$ folds (kinks) has
    \hfill (A) $N-1$ pieces \quad (B) $N$ pieces \quad \textbf{(C) $N+1$ pieces} \ans{$\leftarrow$} \quad (D) $2N$ pieces
  \item Sliding one small filter across the whole input is convolution; reusing that filter everywhere is
    \hfill \textbf{(A) weight sharing} \ans{$\leftarrow$} \quad (B) the bias \quad (C) the loss \quad (D) a rotation
  \item Gradient descent updates a weight by stepping
    \hfill (A) along the slope \quad \textbf{(B) opposite the slope (downhill)} \ans{$\leftarrow$} \quad (C) to a random weight \quad (D) to $0$
  \item A positive covariance $\sigma_{xy}>0$ between two features means they
    \hfill (A) move in opposite directions \quad \textbf{(B) tend to move together} \ans{$\leftarrow$} \quad (C) are unrelated \quad (D) are equal
\end{enumerate}

\begin{teachernote}
Item 1: ReLU zeroes out negatives, so $\relu(-3)=0$. Item 2: a layer is $\relu(A\vv+\bb)$ --- weight, bias, rectify.
Item 3: $N$ folds cut the line into $N+1$ straight pieces. Item 4: reusing one filter everywhere is weight sharing
(why a conv layer has so few weights). Item 5: the slope points uphill, so you step the opposite way (downhill).
Item 6: a positive covariance means the two features rise and fall together.
\end{teachernote}

% ======================================================================
\parthead{Part C --- Short Answer \& Computation (35 pts)}
Show your work.

\begin{enumerate}[leftmargin=1.9em, itemsep=10pt]
  \item One layer $\relu(A\vv+\bb)$.
        \ans{$A\vv=\begin{bsmallmatrix}1&1\\1&-1\end{bsmallmatrix}\begin{bsmallmatrix}1\\3\end{bsmallmatrix}
        =\begin{bsmallmatrix}4\\-2\end{bsmallmatrix}$; add the bias:
        $A\vv+\bb=\begin{bsmallmatrix}4\\-2\end{bsmallmatrix}+\begin{bsmallmatrix}0\\1\end{bsmallmatrix}
        =\begin{bsmallmatrix}4\\-1\end{bsmallmatrix}$; then ReLU keeps positives and zeroes the negative:
        $\relu\begin{bsmallmatrix}4\\-1\end{bsmallmatrix}=\begin{bsmallmatrix}4\\0\end{bsmallmatrix}$.}
  \item Tent function $F(x)=\relu(x)-2\relu(x-2)+\relu(x-4)$.
        \ans{\textbf{(a)} Turn on only the ReLUs whose fold is left of $x$:
        $F(0)=0$, $F(2)=2-0+0=2$, $F(4)=4-2(2)+0=0$, $F(6)=6-2(4)+2=0$.
        \textbf{(b)} Folds at $x=0,2,4$ (where each ReLU switches on); the peak is $(2,2)$ --- the graph goes up,
        then down, then flat (a tent).}
  \item Edge filter $[-1,1]$ across $\vv=[2,2,2,6,6]$.
        \ans{\textbf{(a)} Each output is (next $-$ current): $2{-}2,\ 2{-}2,\ 6{-}2,\ 6{-}6$, so the feature map is
        $[0,0,4,0]$. \textbf{(b)} The one nonzero value ($4$) sits at the jump from $2$ to $6$ --- that is the edge.
        On the flat stretches the two window values are equal, so next $-$ current $=0$.}
  \item Convolution matrix $A=\begin{bsmallmatrix}-1&1&0&0\\0&-1&1&0\\0&0&-1&1\end{bsmallmatrix}$.
        \ans{\textbf{(a)} $A\begin{bsmallmatrix}1\\1\\6\\6\end{bsmallmatrix}
        =\begin{bsmallmatrix}1{-}1\\6{-}1\\6{-}6\end{bsmallmatrix}=\begin{bsmallmatrix}0\\5\\0\end{bsmallmatrix}$ ---
        the same slide-the-filter output. \textbf{(b)} The filter uses only \textbf{2} weights (the $-1$ and $1$,
        reused down every diagonal), while a full $3\times4$ layer would have its own weight in each entry: $12$.}
  \item Gradient descent on $L(w)=(w-4)^2$, $L'(w)=2w-8$, $s=0.25$, $w_0=0$.
        \ans{\textbf{(a)} $w_1=0-0.25(-8)=2$;\ $w_2=2-0.25(-4)=3$;\ $w_3=3-0.25(-2)=3.5$.
        \textbf{(b)} Losses $L(0)=16,\ L(2)=4,\ L(3)=1,\ L(3.5)=0.25$ --- shrinking fast.
        \textbf{(c)} It is heading toward $w=4$, the bottom of the bowl; you know to stop when the slope
        $L'(w)=0$ (no downhill left).}
  \item One gradient step on $L=(w_1-1)^2+(w_2-3)^2$.
        \ans{$\nabla L(0,0)=\big(2(0-1),2(0-3)\big)=(-2,-6)$. Step against it:
        $(0,0)-0.25(-2,-6)=(0.5,1.5)$. The loss bottoms out at $(w_1,w_2)=(1,3)$, where both slopes are $0$.}
  \item Mean, variance, covariance, PCA for $(2,4),(4,2),(6,8),(8,6)$.
        \ans{\textbf{(a)} Mean $=\big(\tfrac{2+4+6+8}{4},\tfrac{4+2+8+6}{4}\big)=(5,5)$.
        \textbf{(b)} Centered $x:-3,-1,1,3$ and $y:-1,-3,3,1$. Variances
        $\sigma_x^2=\tfrac{9+1+1+9}{4}=5$, $\sigma_y^2=5$; covariance
        $\sigma_{xy}=\tfrac{(-3)(-1)+(-1)(-3)+(1)(3)+(3)(1)}{4}=\tfrac{12}{4}=3$.
        \textbf{(c)} $V=\begin{bsmallmatrix}5&3\\3&5\end{bsmallmatrix}$.
        \textbf{(d)} Eigenvalues $\lambda=5\pm3=8,2$; $\mathrm{PC}_1$ (eigenvector $\tfrac{1}{\sqrt2}\begin{bsmallmatrix}1\\1\end{bsmallmatrix}$)
        holds $\dfrac{8}{8+2}=\dfrac{8}{10}=80\%$ of the total variance.}
\end{enumerate}

% ======================================================================
\parthead{Part D --- Extended Response (10 pts)}
Explain your reasoning in full sentences.

\begin{enumerate}[leftmargin=1.9em, itemsep=12pt]
  \item How a network learns.
        \ans{\textbf{(a)} A layer $\relu(A\vv+\bb)$ multiplies the input by weights $A$, shifts by the bias $\bb$
        (which places each fold), then rectifies with ReLU (which zeroes negatives, creating a kink). Each row adds one
        shifted ReLU, so a wider $A$ stacks more folds; a sum of shifted ReLUs is a continuous piecewise-linear function
        whose bends can be positioned and tilted to trace the data. \textbf{(b)} Gradient descent starts from some weights
        and repeatedly steps opposite the slope of the squared-error loss $L$ (the rule $\ww\leftarrow\ww-s\,\nabla L$),
        which lowers $L$ each step --- it rolls downhill in the loss bowl. The bottom of the bowl is where the slope is
        $0$: no downhill direction remains, so those weights give the smallest loss and are the ones the network learns.}
  \item Why the covariance matrix unlocks PCA.
        \ans{\textbf{(a)} $V$ is symmetric because the covariance of $x$ with $y$ equals the covariance of $y$ with $x$
        ($\sigma_{xy}=\sigma_{yx}$), so the off-diagonal entries match and $V=V^{\mathsf T}$. By the spectral theorem
        (Unit 6) a symmetric matrix always has \emph{real} eigenvalues and \emph{perpendicular} eigenvectors.
        \textbf{(b)} Those eigenvectors are the principal components --- the perpendicular directions the data cloud
        spreads along --- and each eigenvalue is the variance (amount of spread) along its direction, largest first.
        The total variance is $\sigma_x^2+\sigma_y^2$, the trace of $V$; since a symmetric matrix's trace equals the sum
        of its eigenvalues, the total spread is split exactly among the principal directions.}
\end{enumerate}

\begin{teachernote}
\textbf{Scoring Part D (5 pts each).} D1: 3 pts (a) layer $=$ weight/bias/ReLU, wider $A$ $=$ more folds $=$ flexible
piecewise-linear fit; 2 pts (b) gradient descent steps against the slope to shrink $L$, slope $0$ $=$ best weights.
D2: 3 pts (a) $\sigma_{xy}=\sigma_{yx}\Rightarrow V$ symmetric $\Rightarrow$ real eigenvalues, perpendicular
eigenvectors; 2 pts (b) eigenvectors $=$ principal components (directions of spread), eigenvalues $=$ variances,
trace $=$ total variance $=$ sum of eigenvalues. Accept equivalent wording throughout.
\end{teachernote}

\end{document}
